\documentclass[UTF8,fontset=windows,11pt,a4paper]{ctexart}
\usepackage{amsmath,amssymb,amsthm,booktabs,geometry,hyperref}
\DeclareFontShape{OMS}{cmsy}{m}{n}{<5-9>cmsy8<9->cmsy10}{}
\geometry{left=27mm,right=27mm,top=25mm,bottom=25mm}
\hypersetup{colorlinks=true,linkcolor=black,citecolor=black,urlcolor=blue}
\setlength{\parskip}{3pt}
\setlength{\emergencystretch}{2em}
\linespread{1.18}
\clubpenalty=10000\widowpenalty=10000
\pagestyle{plain}
\ctexset{section={format=\large\bfseries},subsection={format=\normalsize\bfseries}}
\newtheorem{theorem}{定理}
\renewcommand{\proofname}{证明}
\numberwithin{equation}{section}
\makeatletter
\renewcommand{\@maketitle}{\newpage\null\vskip0em
\begin{center}{\LARGE\heiti\@title\par}\vskip1em
{\normalsize\@date\par}\end{center}\par\vskip0.7em}
\makeatother
\begin{document}
\title{完成可靠学习是否值得投资：证据要求与股东价值}
\author{}\date{2026年10月31日 · 结构与表达重写}
\maketitle
\begin{abstract}
能够支付学习过程，并不意味着完成学习值得投资。本文将一项可靠行动任务与完整股东账目连接：初始资本一次投入，试验须先付款，捕获收入可继续使用，闲置现金按给定利率增长，返款按股东估值率折现。我们区分完成一次付费试验即可作出通过判定，以及通过前必须已有真实捕获两种任务要求。对后一类，用首次捕获前后的现金价值上界证明某个低利率区间内每项有限政策均为负净现值；再将经营收益上界与区分好坏市场所需的信息联立，限制可达的收益规模。有限构造说明某些较高利率下确有正值，但收益保证极小、资源需求巨大。正文集中于可支付、统计可靠和投资值得做三者的区别，不将模型内正值解释成现实盈利机会。
\end{abstract}

\section{研究问题：学得会，是否值得学}
投资者运行一套学习程序，先付试验成本，再观察机会是否捕获。收入可以支付后续试验，余额不足时则无法继续。原融资问题询问至少需要多少启动现金；本文另问，返还给出资人的钱在折现后是否超过最初投入。

这两个问题不能相互替代。程序可能可靠地完成任务，却消耗了过多时间、手续费或资金占用。本文保留逐笔支付约束，新增闲置现金收益和股东估值。它是原融资研究的后续模型；正利息构造不适用于原先闲置现金无收入的制度。

决策者可以选择试验强度、继续或停止、何时返还资金，以及通过后怎样经营。本文研究给定任务内的决策价值，不研究向市场发行股票、债务结构或真实融资报价。

\section{观察、任务与资金}
\subsection{试验与后续经营是两种付款方式}
每次新机会有隐藏捕获速度 $Z\ge0$，来自采购前固定的未知分布 $\nu$，满足 $\mathbb E_\nu Z^2\le M_2$。各新机会独立。试验先选 $q>0$，预付 $cq$；到完整窗口结束，只观察成功与否，成功收款 $1-s$。没有免费观察个体类型、窗口内捕获时刻或额外准确旧读数。

通过后的一次规定经营行动，最多持续到 $A$，捕获即停止支出并收入1，其期望利润为
\begin{equation}
 F_\nu(A)=\mathbb E_\nu(1-e^{-ZA})
 -c\mathbb E_\nu\int_0^A e^{-Zt}\,dt.
\end{equation}
试验费预付且窗口内不退款，与经营捕获后停止支出不同；试验成功费 $s$ 也不能混入经营费用。是否存在这样的两阶段收费差异，需要实际制度依据。

设 $\delta,\Delta>0$，令
\begin{equation}
 \lambda=c+\delta+\Delta,\qquad A=\delta/M_2,
 \qquad g_A=\delta^2/(2M_2).
\end{equation}
好市场满足 $\mathbb E_\nu Z\ge\lambda$；不合格市场由原行动的 $F_\nu(A)<g_A$ 定义。二阶矩界给
\begin{equation}
 F_\nu(A)\ge(\mathbb E_\nu Z-c)A-M_2A^2/2,
\end{equation}
因此好市场确有行动利润余量，但这不是完整学习项目的利润。两类之间允许间隔。对每个不合格市场，通过概率须不超过 $\alpha$；对每个好市场，拒绝概率须不超过 $\beta$。

\subsection{为什么区分是否必须真的捕获过}
第一类任务要求通过前至少完成一次付费正强度试验，随后可以根据预定规则判断；不要求日志中已经成功。第二类额外要求，通过前必须至少有一次完成的真实付费捕获。第二类是更强的任务条件，第一类不能被事后按第二类排除。

这一差别不是文字上的证明格式。若允许在零成功记录上通过，政策可在错误预算内采用少量提前通过的随机分支；要求真实成功会禁止这种做法。是否应禁止，须由实际验收目的决定。本文比较两类，不默认真实捕获要求普遍必要。

\subsection{股东究竟投入和收回什么}
股本 $B$ 在时刻0一次缴清，不追加、不借款。现金按连续收益率 $r_C$ 增长，股东按 $\rho\ge r_C$ 折现。若在时刻 $t_j$ 返还 $D_j$，则政策的最坏好市场净现值为
\begin{equation}\label{eq:npv}
 V_\pi=\inf_{\nu\ {
m good}}
 \mathbb E_\nu\left[\sum_j e^{-\rho t_j}D_j\right]-B.
\end{equation}
返款包括最终清算，不能超过当时现金。每项政策的资本、机会数和最终期限均有限；可以比较越来越长但各自有限的政策。通过时须持有规定行动的预付金 $cA$。后续可经营或退出，不强迫重复完成原行动。

正值指式\eqref{eq:npv}大于零，不是账户最后有正余额，也不是某条成功路径赚钱。只将现金存放到 $T$ 后取回，净现值是 $B[e^{-(\rho-r_C)T}-1]\le0$，所以正息本身不等于免费投资利润。

\section{为什么取得证据可能不值得}
\subsection{单次试验的损失还不是完整证明}
固定一个合法好市场 $G$：$Z$ 恒等于 $\lambda$。短试验每单位曝光的期望净收入不超过 $(1-s)\lambda-c$。当前参数 $c=1.5,\lambda=1.505,s=0.01$，该值为 $-0.01005$；但通过后的无成功费经营具有正的收入率余量 $\lambda-c=0.005$。

因此，业务本身可能有利，而学习阶段仍消耗价值。仅凭一次试验亏损，还不能排除反馈、返款和后续经营补偿它。下面用同一个市场的整个现金过程证明上界。

\subsection{用可核对的价值上界覆盖全部政策}
为证明不可能性，先给投资者更有利的条件：成本可连续支付，捕获即时揭示并回款，首次学习成功后免除剩余检查要求直接进入经营。放宽后的项目都不能获利时，原任务也不能。生成式与返款导数不等式属于已有分红控制验证方法\cite{dividend}；本文还须单独处理首次学习成功、现金利息和可靠任务，不能直接套用既有障碍最优性结论。

令 $x$ 为当前现金，用 $x+h(x)$ 作为首次捕获后的折现总回款上界，其中 $h\ge0$ 非减且 $h(0)=0$；首次捕获前只用现金 $x$。$h$ 是证明工具，表示给后续经营的额外价值额度，不是额外现金。令
\begin{equation}
 d_s=c/\lambda-(1-s)>0,\qquad\zeta=\rho-r_C.
\end{equation}

\begin{theorem}[真实捕获要求下的负净现值判据]\label{thm:negative}
假设非负非减、连续且局部绝对连续的分段光滑 $h$ 在整个 $x\ge0$ 及拐点两侧满足
\begin{align}
 &(r_Cx-c)(1+h'(x))
 +\lambda[1+h(x+1)-h(x)]-\rho[x+h(x)]\le0,\label{eq:operate}\\
 &r_Cx(1+h'(x))-\rho[x+h(x)]\le0,\label{eq:idle}\\
 &\lambda[h(x+1-s)-d_s]-\zeta x\le-\epsilon_0,
 \qquad\epsilon_0>0.\label{eq:learn}
\end{align}
则每项有限且通过前必须真实捕获的可靠政策都有 $V_\pi<0$。若检查在 $T$ 前结束，则
\begin{equation}
 V_\pi\le-\frac{\epsilon_0(1-\beta)}\lambda e^{-\rho T}<0.
\end{equation}
\end{theorem}

\begin{proof}
式\eqref{eq:operate}和\eqref{eq:idle}分别说，首次捕获后的经营、闲置不会提高折现价值上界的期望。返款 $d$ 会把现金从 $x$ 减到 $x-d$，因 $1+h'\ge1$，返款加剩余价值也不增加上界。

首次捕获前，上界为 $x$。在学习期间，现金连续变化率为 $r_Cx-c$；以强度 $\lambda$ 成功时，现金增加 $1-s$ 并切换到 $x+1-s+h(x+1-s)$。扣除股东折现后的期望变化率恰为式\eqref{eq:learn}左侧，至多 $-\epsilon_0$。不学习时只产生非正利差损失。

对折现价值加累计返款求期望，得到相对于初始 $B$ 的损失至少为首次成功前学习时间的折现积分乘 $\epsilon_0$。成功强度为 $\lambda$，其补偿恒等式将该积分期望写成首次成功折现概率除以 $\lambda$。可靠性与真实捕获要求共同保证截至 $T$ 的成功概率至少 $1-\beta$，故该折现概率至少 $(1-\beta)e^{-\rho T}$。有限期限和容量保证可先在有界过程上使用期望恒等式；分段点由左右导数和实际跳变量处理。最后，$G$ 为好市场之一，最坏好侧值不超过 $G$ 的值。
\end{proof}

\noindent\textbf{投资含义。}
这是对全部允许反馈政策的上界，而非只说明某个学习器不划算。当前有理折线 $h$ 在无限现金半轴上的核查给出：$0\le r_C/\rho\le0.08225$ 时，每项有限真实捕获政策均为负值。这里0.08225是两种连续利率的比例，不是年利率8.225\%。

若 $r_C>0$，可以用任意小初始存款等待增长，再支付完整学习任务；返还资金非负，净现值可从负侧逼近零。因此这个区间内上确界为零却没有合格有限政策达到零。它不是正收益方案；若允许一开始不承担任务，退出本来就给零值。

\section{可靠性还限制能赚到多少}
仅证明存在极小正值，不足以说明能达到给定投资目标。记全部好市场要求的完整目标为 $G_0=0.1$；它不同于定义单次行动合格的 $g_A$。

仍在同一 $G$，设后续经营的额外价值上界为 $b_*\ge0$。令 $\tau$ 为通过或退出的时刻，$J$ 为通过指示，定义
\begin{equation}
 D_\rho=\mathbb E_G[e^{-\rho\tau}J],\qquad
 C_\rho=\mathbb E_G\sum_i e^{-\rho T_i}X_i,
\end{equation}
其中 $X_i$ 是第 $i$ 个完成学习窗口是否成功，$T_i$ 是结果实际揭示和回款时间。$D_\rho$ 是折现后的通过权重，不是普通通过概率；$C_\rho$ 是折现成功次数，不是费用金额。

现金比较给净现值上界 $b_*D_\rho-d_sC_\rho$。再选固定不合格市场 $Q$：$Z\equiv\ell=c+\delta/2$。虽然 $\ell>c$，其规定行动利润仍小于 $g_A$，所以确属原不合格类，不能仅凭均值是否超过成本分类。

每项 $G$ 成功所能支持的有利信息上界为
\begin{equation}
 I_* =\log(\lambda/\ell)-1+\ell/\lambda.
\end{equation}
引入独立、均值为 $1/\rho$ 的指数截止，将折现权重写成截止前事件概率。完整日志的似然比较和事件合并给
\begin{equation}\label{eq:info}
 \operatorname{kl}_+(D_\rho,\alpha)\le I_*C_\rho.
\end{equation}
这里 $\operatorname{kl}(a,b)=a\log(a/b)+(1-a)\log[(1-a)/(1-b)]$；当 $a\le\alpha$ 时，$\operatorname{kl}_+(a,\alpha)$ 定义为零，否则取二元相对熵。指数截止只是证明工具，没有给真实投资者免费消息。

真实捕获要求另给逐路径关系 $C_\rho\ge D_\rho$，因为通过前至少有一项较早完成的成功，其折现权重不小于通过权重。因此
\begin{equation}
 V_\pi\le\sup_{0\le D\le1}
 \left\{b_*D-d_s\max\left[D,\frac{\operatorname{kl}_+(D,\alpha)}{I_*}\right]\right\}.
\end{equation}
不要求真实捕获的任务只能删掉方括号中的 $D$，不能使用同一加强界。

\noindent\textbf{为什么这项连接重要。}
资本增长可以改变何时付得起费用，却不能免费增加区分好坏市场的信息。现金损失和信息要求必须在同一个市场中连接，不能把不同最坏市场的期望费用相加。这里的经营上界和信息上界并不声称精确最优，但足以排除部分收益目标。

\section{有限构造说明什么，不说明什么}
当前参数为 $M_2=25,\rho=0.00249,\alpha=\beta=0.05$，其余如上。正文只保留与投资判断直接有关的结果：
\begin{center}
\begin{tabular}{p{0.29\linewidth}p{0.61\linewidth}}\toprule
问题 & 现有结论及限制\\\midrule
真实捕获任务能否正值 & 低比例不超过0.08225时每项有限政策负值；比例至少0.0875时有正值构造；中间临界点未定。\\
能否达到 $G_0=0.1$ & 比例不超过0.9998时，两类任务均被现金与信息上界排除。\\
不要求真实捕获的任务 & 比例至少0.9999995时有完整有限政策保证大于0.1000029。\\
要求真实捕获的任务 & 比例至少0.9999999时现有保证大于0.09999105，仍未达到0.1。\\\bottomrule
\end{tabular}
\end{center}

较低正值构造的保证仅约 $10^{-15}$，不能解释为有实际投资吸引力。接近0.1的构造采用约603初始储备、数千亿机会和极长模型期限。普通任务的早期规则还使用接近允许错误上限的小概率通过分支：数学上属于该任务，但是否对应合理的业务验收需要另议。有限资源不等于现实资源可得。

本次在隔离副本实际重跑低利率负值函数、0.0875比例正值构造、0.1目标必要界、普通任务有限储备政策和真实捕获近目标政策的核验程序；现金转移、折现账目和好坏侧检验界通过。这是对既有完整数值对象的重算与核查，不是重新进行全部候选搜索。一般性证明仍须审读；执行日志、输入散列和退出状态随稿提供。

\section{与原问题的关系及剩余工作}
原文研究无安全现金收入时，为学习和行动至少准备多少资金；本稿增加现金收益率，转向完整股东价值。两者不能互相替代。现金收益率、学习成功费、真实捕获要求以及允许任意长有限等待，都是会影响结论的实质条件。

时间与机会容量的联合信息限制、极高利率下的具体储备构造和更多政策变体移至附件。正文保留两条可以逐步复述的论证：真实捕获前的损失是否能被后续经营补偿；完成可靠判断所需的信息是否已经消耗了可取得的价值。现实解释应先检验这些制度，而不是先以极小正值宣称投资可行。

\appendix
\section{折现信息界为什么与成功次数相连}
固定正文 $G:Z=\lambda$ 和 $Q:Z=\ell$。先给学习者更丰富的记录：每窗除了是否成功，还看到截止 $q$ 的首次捕获时间。若 $T$ 是捕获时间，该记录由成功密度与未成功概率组成。在 $G$ 中其相对 $Q$ 的相对熵为
\begin{align}
 K_q&=(1-e^{-\lambda q})\log(\lambda/\ell)
 -(\lambda-\ell)\mathbb E_G\min(T,q)\\
 &=\left[\log(\lambda/\ell)-1+\ell/\lambda\right](1-e^{-\lambda q}).
\end{align}
真实程序只看到完整窗口的二元结果，信息不超过此式。该放宽只用于必要界，没有改变真实数据权限。

取独立 $E\sim\operatorname{Exp}(\rho)$，只保留完成时刻 $T_i\le E$ 的原日志。给定过去和当前选择，窗口完成时刻在当前结果揭示前确定，且 $\Pr(E\ge T_i)=e^{-\rho T_i}$。链式相对熵将各完成窗口的信息相加，得到总信息至多 $I_*C_\rho$。在截止前通过的事件概率，在 $G$ 为 $D_\rho$，在 $Q$ 不超过原坏侧上限 $\alpha$；数据处理与二元相对熵单调性给正文式\eqref{eq:info}。若 $D_\rho\le\alpha$，仅使用零信息下界。

\section{现金上界与信息上界必须在同一市场中连接}
在 $G$ 中，每次成功概率 $p_G(q)\le\lambda q$，预付费用在结果完成前发生，其折现成本至少为完成时刻折现的 $cq$。平均成本因此至少是每次成功 $c/\lambda$；成功净收 $1-s$，每次成功平均至少消耗 $d_s$ 的折现现金。由于 $r_C\le\rho$，等待利息不会提高现金的股东现值。

若通过后的经营回收至多是当时现金加额外价值 $b_*$，累计返款和清算的同一现金账目给
\begin{equation}
 V_\pi\le b_*D_\rho-d_s C_\rho.
\end{equation}
实际捕获要求还给 $C_\rho\ge D_\rho$；不要求捕获时不能使用它。令 $K_*=d_s/I_*$，不要求捕获的上界可显式优化为
\begin{equation}
 \sup_{0\le D\le1}\{b_*D-K_*\operatorname{kl}_+(D,\alpha)\}
 =K_*\log(1-\alpha+\alpha e^{b_*/K_*}).
\end{equation}
这是对 $D$ 求导并解其对数比值所得。$b_*\ge0$ 时最优点不小于 $\alpha$，故截为零的分支不改变最大值。加强后的真实捕获式保留一维上确界，并不假装已经显式解出。

一个统一但较松的后续比较是 $x+(\lambda-c)/\rho$。在更有利的连续经营模型中，其经营生成式扣折现恰为 $-(\rho-r_C)x\le0$，闲置也非正，返款导数等于1；因此可用于上界。它在零现金额外赠送正值，不能拿来替代正文定理\ref{thm:negative}所需的 $h(0)=0$ 严格损失函数。

\section{数值表的证明责任与完整资源}
正文各结论的原完整政策和函数不等式，分别保存在技术材料中。负值方向覆盖整个现金半轴，不止网格点；正值方向把有限尾部、每笔付款和后段学习纳入同一初始股本。

超过0.1的普通任务构造初始储备约603，最迟模型时刻40,015,000，最多250,000,000,001次机会。真实捕获近目标构造初始股本603.0512606665，最迟时刻80,015,000，最大容量500,034,173,771。模型时刻没有自动对应天或年，不能自行加日历单位。

本次保留的数值是可行保证或不可能性界，不是实际最优值。构造中的高利率、极长等待和允许的早期通过规则都应在经济解释中明示。候选搜索结果、现有有理核验和一般性证明分别保存，不能用再次运行同一检查器代替同行评议。

\begin{thebibliography}{9}
\bibitem{dividend} Bayraktar E., Kyprianou A. E., Yamazaki K. On Optimal Dividends in the Dual Model. 2013. 本文仅使用其验证方法作为直接背景；\S2.3式(2.16)、定理2.1的条件与本模型不同。\href{https://warwick.ac.uk/fac/sci/statistics/staff/academic-research/kyprianou/astin.pdf}{作者原文}。
\end{thebibliography}
\end{document}
